\begin{abstract}
Deep learning for infrared meibography segmentation has advanced rapidly, but the comparability of published results is limited by three neglected issues: image-level splits that let the same patient appear in training and test, unreported tiling parameters at inference, and a gap between pixel metrics and clinical Meiboscore. We establish a patient-independent, reproducible benchmark on the public MGD-1K dataset: \nTotal{} deduplicated pairs from 338 subjects, split by patient into \nTrain/\nVal/\nTest{} (\nTrainPatients/\nValPatients/\nTestPatients{} patients), with the split file and assignment signature released.

We first quantify a hidden leakage: the official \officialTestImages-case Meiboscore test list contains \officialTestPatients{} patients, while \officialLeakImages{} of the remaining 800 images come from the same patients; even a patient-level training set still contains \officialLeakPatients/\officialTestPatients{} of the test patients. We then compare UNet, Attention U-Net, WRCANet, MGD-PriorNet, and MGD-Former under one preprocessing, one patient-level split, one loss, and one inference protocol; WRCANet did not converge within the 8\,GB budget, so the main table ranks the four comparable routes and reports WRCANet separately as a feasibility counterexample. MGD-Former uses an ADE20K-pretrained MiT-B2 encoder and a stride-4 high-resolution fusion decoder.

We treat tile size as a first-class experimental variable: with identical weights, changing only tile height moves Dice from 0.370 to 0.655. This motivates Multi-Tile Ensemble (MTE), which averages probabilities from 256/384/512 tiles. MTE improved HD95 but not Dice. We further compare multi-tile variance, predictive entropy, and 8$\times$ TTA by AUROC for failure detection, and report agreement between segmentation-derived R\_loss/Meiboscore and six-round expert grading, together with upper/lower eyelid subgroup differences. Clinical agreement is only QWK \qwkMF, showing that pixel metrics cannot replace expert grading; TTA8 AUROC is near 1.0 but suffers from label--signal coupling and is not direct triage evidence. All accuracy values come from unified training and a patient-level independent test; no unmeasured performance is reported.

\vspace{0.6em}
\noindent\textbf{Keywords:} meibomian gland segmentation; public benchmark; patient-level leakage; tiling ensemble; uncertainty; clinical grading
\end{abstract}
